\section{ANALISIS KOMPETITOR}

\subsection{Analisis Kompetitor}

Berikut adalah perbandingan Savora dengan platform serupa yang beroperasi di Indonesia berdasarkan data aktual per 2025--2026 \cite{mobilefoodwaste2024review}\cite{kusumawati2025mobile}:

\begin{longtable}{|>{\RaggedRight}p{0.22\textwidth}|>{\RaggedRight}p{0.26\textwidth}|>{\RaggedRight}p{0.26\textwidth}|>{\RaggedRight}p{0.26\textwidth}|}
\caption{Perbandingan Fitur dengan Kompetitor}
\label{tab:kompetitor} \\
\hline
\textbf{Aspek / Fitur}
  & \textbf{Savora}
  & \textbf{Surplus Indonesia}
  & \textbf{Garda Pangan} \\
\hline
\endfirsthead
\multicolumn{4}{l}{\small\textit{(Tabel \ref{tab:kompetitor} dilanjutkan)}} \\
\hline
\textbf{Aspek / Fitur}
  & \textbf{Savora}
  & \textbf{Surplus Indonesia}
  & \textbf{Garda Pangan} \\
\hline
\endhead
\hline
\multicolumn{4}{r}{\small\textit{(bersambung ke halaman berikutnya)}} \\
\endfoot
\hline
\endlastfoot
Jenis platform
  & Marketplace food rescue UMKM (MVP)
  & Clearance sale app (recommerce)
  & Food bank sosial / nirlaba \\
\hline
Tersedia di Indonesia
  & \checkmark\ (MVP)
  & \checkmark\ (aktif)
  & \checkmark\ (Surabaya \& sekitar) \\
\hline
Fokus UMKM lokal kuliner
  & \checkmark
  & \checkmark\ (UMKM, horeka)
  & \checkmark\ (distribusi gratis) \\
\hline
Transparansi produk
  & Tinggi (foto, deskripsi, kelayakan)
  & Sedang (nama \& foto)
  & Tinggi (SOP seleksi ketat) \\
\hline
Food Trust Index / info kelayakan
  & \checkmark\ (rule-based, tampil ke user)
  & \ding{55}
  & Seleksi internal (tidak ke user) \\
\hline
Metode pembayaran
  & Cashless via Midtrans sandbox + service fee 5\%
  & Digital saja (OVO, GoPay, DANA)
  & Tidak ada (donasi) \\
\hline
Klasifikasi keyword ulasan
  & \checkmark\ (Aman/Warning/Gawat)
  & \ding{55}
  & \ding{55} \\
\hline
Dynamic discount berbasis kelayakan
  & \checkmark\ (rule-based, 10--50\%)
  & \checkmark\ (hingga 80\%, tidak berbasis kelayakan)
  & \ding{55} \\
\hline
Pickup code \& batas waktu
  & \checkmark\ (pickup code + deadline)
  & Window waktu
  & \ding{55} \\
\hline
Waste Log untuk mitra UMKM
  & \checkmark
  & \ding{55}
  & \ding{55} \\
\hline
Impact tracking ke pengguna
  & \checkmark\ (estimatif, bukan klaim resmi)
  & \checkmark\ (estimasi CO\textsubscript{2}/CH\textsubscript{4})
  & \ding{55} \\
\hline
Rating \& ulasan
  & \checkmark
  & \checkmark
  & \ding{55} \\
\hline
Skala (2025--2026)
  & MVP lomba
  & Ratusan ribu user, Jawa/Bali/Sulsel
  & Lokal Surabaya \\
\hline
\end{longtable}

\subsection{Analisis SWOT}

\begin{table}[htbp]
\centering
\caption{Analisis SWOT Savora}
\label{tab:swot}
\begin{tabular}{|>{\RaggedRight\arraybackslash}p{0.46\textwidth}|>{\RaggedRight\arraybackslash}p{0.46\textwidth}|}
\hline
\textbf{Strengths (Kekuatan)} & \textbf{Weaknesses (Kelemahan)} \\
\hline
\begin{minipage}[t]{\linewidth}
\begin{itemize}[leftmargin=*,nosep]
    \item Fokus masalah jelas dan lokal (UMKM kuliner)
    \item Food Trust Index transparan dan mudah dijelaskan
    \item Food Score Decay memberikan urgensi visual \textit{real-time}
    \item Keyword classification dari ulasan (\textit{crowdsourced safety})
    \item Model pendapatan jelas: \textit{service fee} 5\% + iklan
    \item Cashless via Midtrans sandbox menunjukkan kesiapan integrasi
\end{itemize}
\end{minipage} &
\begin{minipage}[t]{\linewidth}
\begin{itemize}[leftmargin=*,nosep]
    \item Bergantung pada kejujuran input data UMKM
    \item Integrasi Midtrans masih \textit{sandbox}
    \item Validasi data \textit{real} terbatas pada demo/pilot
    \item Butuh waktu untuk adopsi awal
\end{itemize}
\end{minipage} \\
\hline
\textbf{Opportunities (Peluang)} & \textbf{Threats (Ancaman)} \\
\hline
\begin{minipage}[t]{\linewidth}
\begin{itemize}[leftmargin=*,nosep]
    \item Kesadaran \textit{food waste} dan belanja hemat meningkat
    \item Banyak UMKM belum punya kanal khusus makanan surplus
    \item Dapat dikembangkan menjadi PWA tanpa instalasi \textit{native}
    \item Mitra donasi membuka peluang kerja sama sosial
\end{itemize}
\end{minipage} &
\begin{minipage}[t]{\linewidth}
\begin{itemize}[leftmargin=*,nosep]
    \item Marketplace besar dapat menambahkan fitur serupa
    \item \textit{Trust issue} terhadap makanan surplus menghambat adopsi
    \item Customer \textit{no-show} dapat merugikan UMKM
\end{itemize}
\end{minipage} \\
\hline
\end{tabular}
\end{table}

\subsection{Model Bisnis}

Savora mengadopsi model bisnis yang realistis dan dapat didemonstrasikan pada MVP, dengan dua
sumber pendapatan utama:

\begin{enumerate}[leftmargin=*]
    \item \textbf{Service Fee 5\%} : Biaya layanan 5\% yang otomatis ditambahkan pada setiap transaksi
    produk dan ditampilkan transparan di halaman checkout. Total pendapatan \textit{service fee}
    tercatat pada dashboard keuangan admin.
    \item \textbf{Pendapatan Iklan} : Slot iklan \textit{premium listing} untuk UMKM dan iklan pihak
    ketiga (brand luar), yang seluruhnya melalui \textit{approval} admin sebelum tayang.
\end{enumerate}

Model ini sengaja dijaga sederhana agar dapat diuji dalam durasi lomba tanpa \textit{overclaim}.
Sumber pendapatan lanjutan seperti langganan premium UMKM, program CSR korporat, dan kerja sama
kelembagaan diposisikan sebagai \textit{roadmap} setelah MVP tervalidasi.

\subsection{Rencana Keberlanjutan}

\begin{enumerate}[leftmargin=*]
    \item \textbf{Tahap MVP (Lomba)} : Membuktikan kelayakan alur \textit{end-to-end} (listing, cashless
    Midtrans sandbox, pickup, review) dengan data demo/pilot yang realistis.

    \item \textbf{Tahun Pertama} : Fokus akuisisi UMKM dan pengguna di satu kota, mengaktifkan mode
    \textit{production} Midtrans, dan mencapai \textit{break even} dari \textit{service fee} + iklan.

    \item \textbf{Tahun Kedua} : Ekspansi ke beberapa kota besar, memperkuat mitra donasi, serta
    menambah paket langganan premium dan program CSR.

    \item \textbf{Jangka Panjang} : Menjadi platform referensi \textit{food rescue} UMKM di Indonesia
    dengan potensi pengembangan produk turunan (pelatihan UMKM, konsultasi \textit{sustainability}).
\end{enumerate}
